%% Title and Preprocessing
\documentclass[titlepage]{article}\pagestyle{empty}
\usepackage{amsmath}

\author{Kevin Fang}
\title{Homework \#19}
\date{\today}
\usepackage[left=1cm, right=2cm, vmargin=1cm]{geometry}
\usepackage{amsfonts}

\def\therefore{\boldsymbol{\text{ }
\leavevmode
\lower0.4ex\hbox{$\cdot$}
\kern-.5em\raise0.7ex\hbox{$\cdot$}
\kern-0.55em\lower0.4ex\hbox{$\cdot$}
\thinspace\text{ }}}

\usepackage{mathtools}
\DeclarePairedDelimiter\ceil{\lceil}{\rceil}
\DeclarePairedDelimiter\floor{\lfloor}{\rfloor}

% Start Document Here
\begin{document}
\maketitle

\pagebreak
\section*{Question 1.}
The UNIX $fork()$ function is a fundamental system call used in network programming and process management. It is designed to create a new process by duplicating the existing process, known as the parent process. This new process is referred to as the child process. The $fork()$ function is essential in environments where threading support is limited or unavailable, as it allows for concurrent execution of processes.

When a program calls $fork()$, the operating system performs several key actions:

\begin{enumerate}
  \item Process Duplication: The $fork()$ call results in the creation of a child process that is an exact copy of the parent process. This includes duplicating the process's memory space, file descriptors, and execution context.


  \item Copy-on-Write Mechanism: To optimize performance and resource usage, UNIX systems often use a "copy-on-write" strategy. This means that the memory pages of the parent process are not immediately copied for the child. Instead, both processes share the same memory pages marked as read-only. Actual copying occurs only when either process attempts to modify a shared page.


  \item Separate Process IDs: The child process receives a unique process identifier (PID) distinct from the parent. This PID is used to manage and control the child process independently.
\end{enumerate}

Process States After Fork

After invoking $fork()$, both the parent and child processes enter specific states within the operating system:
\begin{enumerate}
  \item Running State: Immediately after the fork, both processes are in the running state. They execute the next instruction following the $fork()$ call. The return value of $fork()$ helps distinguish between the parent and child:
    \begin{itemize}
      \item The parent receives the child's PID as the return value.
      \item The child receives a return value of zero.
    \end{itemize}
  \item Ready State: If the CPU is not available, either process may enter the ready state, waiting for CPU time to continue execution.
  \item Blocked State: Processes can enter a blocked state if they are waiting for an event, such as I/O completion. This state is independent of the $fork()$ operation but is common in process lifecycle management.
  \item Terminated State: Eventually, processes will complete their execution and enter the terminated state. The parent process can use system calls like $wait()$ to retrieve the termination status of the child process, ensuring proper cleanup.
\end{enumerate}
Windows does not implement the UNIX-style $fork()$ system call. The concept and mechanism of $fork()$—duplicating a running process, including its memory and execution context—are specific to UNIX and UNIX-like operating systems (such as Linux and macOS).
Windows uses a different model for creating new processes:

\begin{enumerate}
  \item CreateProcess API:
The primary way to create a new process in Windows is via the $CreateProcess()$ function. This API starts a new process from an executable file, rather than duplicating the current process.


  \item No Memory Duplication:
Unlike $fork()$, $CreateProcess()$ does not copy the parent’s memory space, file descriptors, or execution state. Instead, it loads a fresh instance of the specified program.


  \item Explicit Inheritance:
If you want the child process to inherit certain handles or environment variables, you must specify this explicitly when calling $CreateProcess()$.
\end{enumerate}

\begin{enumerate}
  \item UNIX ($fork()$):

    \begin{itemize}
      \item Duplicates the current process.
      \item Both parent and child continue executing from the next instruction.
      \item Useful for servers that spawn child processes to handle requests.
    \end{itemize}
  \item Windows ($CreateProcess()$):

    \begin{itemize}
      \item Starts a new process from a given executable.
      \item Parent and child do not share memory or execution context unless explicitly set up.
      \item Concurrency Models:
Windows encourages the use of threads (via $CreateThread()$) and asynchronous I/O for concurrency, rather than process duplication.
    \end{itemize}
\end{enumerate}
\end{document}